\fichatitulo{16 · APLICAÇÃO}{Preprint 338 · 2025}{The World Avatar no parlamento}
{\small Mosbach et al. \textit{Parliamentary debates in The World Avatar: A hybrid retrieval-augmented generation system}. Preprint 338 · 2025. \href{https://como.ceb.cam.ac.uk/media/preprints/c4e-preprint-338.pdf}{Fonte consultada}.\par}

\campo{Problema e objetivo}
Como responder a perguntas sobre discursos do Bundestag e suas relações?

\campo{Principal contribuição · síntese interpretativa}
Combina busca vetorial sobre discursos e grafo de conhecimento para explorar conteúdo e metadados parlamentares.

\campo{Método / natureza da evidência}
Prova de conceito que compara abordagens de consulta e examina a adequação das respostas; inclui experiências informais com usuários.

\campo{Resultado ou argumento principal}
A discussão mostra que perguntas relacionais e de contagem exigem tratamento diferente de perguntas respondidas por trechos individuais.

\campo{Excertos-chave · seleção literal}
\begin{itemize}
\excerto{1}{our informal experiments with human users}{Discussão.}
\excerto{2}{it should not be concluded that this is necessarily due to the use of knowledge graphs}{Discussão.}
\end{itemize}

\campo{Limitações e análise crítica}
Análise da ficha: métricas que valorizam referências a discursos podem desfavorecer respostas agregadas corretas. Alterações de arquitetura e de prompts dificultam atribuir diferenças exclusivamente ao grafo.

\campo{Aproveitamento na dissertação · proposta}
Trabalhos relacionados: separar perguntas factuais, agregações e relações entre entidades.

\registro{Fonte consultada nesta preparação: Resumo e discussão/conclusão; consulta parcial. Chave: \texttt{mosbachEtAl2025worldAvatarParliament}.}
